\documentclass[11pt]{article}
\usepackage{amsmath, amssymb, amsthm}
\usepackage[margin=1in]{geometry}
\usepackage{enumitem}

\newcommand{\bA}{\mathbf{A}}
\newcommand{\bB}{\mathbf{B}}
\newcommand{\bC}{\mathbf{C}}
\newcommand{\bX}{\mathbf{X}}
\newcommand{\bSigma}{\boldsymbol{\Sigma}}
\newcommand{\bmu}{\boldsymbol{\mu}}

\begin{document}
\section*{Problem 2.2}

Given
\[
\bA = \begin{bmatrix} -1 & 3 \\ 4 & 2 \end{bmatrix}, \qquad
\bB = \begin{bmatrix} 4 & -3 \\ 1 & -2 \\ -2 & 0 \end{bmatrix}, \qquad
\bC = \begin{bmatrix} 5 \\ -4 \\ 2 \end{bmatrix}.
\]

\paragraph{(a) $5\bA$}
\[
5\bA = \begin{bmatrix} -5 & 15 \\ 20 & 10 \end{bmatrix}.
\]

\paragraph{(b) $\bB\bA$}
$\bB$ is $3\times 2$ and $\bA$ is $2\times 2$, so $\bB\bA$ is $3\times 2$:
\[
\bB\bA =
\begin{bmatrix}
4(-1)+(-3)(4) & 4(3)+(-3)(2) \\
1(-1)+(-2)(4) & 1(3)+(-2)(2) \\
-2(-1)+0(4) & -2(3)+0(2)
\end{bmatrix}
=
\begin{bmatrix}
-16 & 6 \\
-9 & -1 \\
2 & -6
\end{bmatrix}.
\]

\paragraph{(c) $\bA'\bB'$}
Since $\bA'\bB' = (\bB\bA)'$,
\[
\bA'\bB' = \begin{bmatrix} -16 & -9 & 2 \\ 6 & -1 & -6 \end{bmatrix}.
\]

\paragraph{(d) $\bC'\bB$}
$\bC' = [5,\,-4,\,2]$ is $1\times 3$, $\bB$ is $3\times2$, so $\bC'\bB$ is $1\times 2$:
\[
\bC'\bB = \big[\,5(4)+(-4)(1)+2(-2),\;\; 5(-3)+(-4)(-2)+2(0)\,\big] = [\,12,\, -7\,].
\]

\paragraph{(e) Is $\bA\bB$ defined?}
$\bA$ is $2\times 2$ and $\bB$ is $3\times 2$. For the product $\bA\bB$ we would need the number of columns of $\bA$ (which is 2) to equal the number of rows of $\bB$ (which is 3). Since $2 \neq 3$, \textbf{$\bA\bB$ is not defined.}

\section*{Problem 2.4}

\paragraph{(a) $(\bA')^{-1} = (\bA^{-1})'$}
Start from $\bA\bA^{-1} = \bI$. Taking transposes of both sides and using $\bI' = \bI$,
\[
(\bA\bA^{-1})' = \bI' \;\Longrightarrow\; (\bA^{-1})'\bA' = \bI.
\]
Similarly, from $\bA^{-1}\bA = \bI$,
\[
(\bA^{-1}\bA)' = \bI \;\Longrightarrow\; \bA'(\bA^{-1})' = \bI.
\]
Thus $(\bA^{-1})'$ acts as both a left and a right inverse of $\bA'$. Since the inverse of a matrix is unique,
\[
(\bA')^{-1} = (\bA^{-1})'.
\]

\paragraph{(b) $(\bA\bB)^{-1} = \bB^{-1}\bA^{-1}$}
Compute
\[
(\bB^{-1}\bA^{-1})(\bA\bB) = \bB^{-1}(\bA^{-1}\bA)\bB = \bB^{-1}\bI\bB = \bB^{-1}\bB = \bI,
\]
\[
(\bA\bB)(\bB^{-1}\bA^{-1}) = \bA(\bB\bB^{-1})\bA^{-1} = \bA\bI\bA^{-1} = \bA\bA^{-1} = \bI.
\]
Since $\bB^{-1}\bA^{-1}$ is both a left and right inverse of $\bA\bB$, and inverses are unique,
\[
(\bA\bB)^{-1} = \bB^{-1}\bA^{-1}.
\]

\section*{Problem 2.8}

Given $\bA = \begin{bmatrix} 1 & 2 \\ 2 & -2 \end{bmatrix}$.

\paragraph{Eigenvalues}
\[
|\bA - \lambda \bI| = (1-\lambda)(-2-\lambda) - 4 = \lambda^2 + \lambda - 6 = 0
\;\Longrightarrow\; (\lambda-2)(\lambda+3) = 0.
\]
\[
\lambda_1 = 2, \qquad \lambda_2 = -3.
\]

\paragraph{Eigenvectors}
For $\lambda_1 = 2$: solve $(\bA - 2\bI)v = 0$:
\[
\begin{bmatrix} -1 & 2 \\ 2 & -4 \end{bmatrix} v = 0 \;\Longrightarrow\; v_1 = 2v_2.
\]
Taking $v = (2,1)'$ and normalizing ($\|v\| = \sqrt{5}$):
\[
e_1 = \left(\frac{2}{\sqrt5}, \frac{1}{\sqrt5}\right)'.
\]

For $\lambda_2 = -3$: solve $(\bA + 3\bI)v = 0$:
\[
\begin{bmatrix} 4 & 2 \\ 2 & 1 \end{bmatrix} v = 0 \;\Longrightarrow\; v_2 = -2v_1.
\]
Taking $v = (1,-2)'$ and normalizing ($\|v\| = \sqrt5$):
\[
e_2 = \left(\frac{1}{\sqrt5}, \frac{-2}{\sqrt5}\right)'.
\]

\paragraph{Spectral decomposition}
\[
\bA = \lambda_1 e_1 e_1' + \lambda_2 e_2 e_2'
= 2\cdot\frac{1}{5}\begin{bmatrix} 4 & 2 \\ 2 & 1 \end{bmatrix}
+ (-3)\cdot\frac{1}{5}\begin{bmatrix} 1 & -2 \\ -2 & 4 \end{bmatrix}.
\]
Numerically,
\[
\bA = \frac{1}{5}\begin{bmatrix} 8 & 4 \\ 4 & 2 \end{bmatrix} + \frac{1}{5}\begin{bmatrix} -3 & 6 \\ 6 & -12 \end{bmatrix}
= \frac{1}{5}\begin{bmatrix} 5 & 10 \\ 10 & -10 \end{bmatrix} = \begin{bmatrix} 1 & 2 \\ 2 & -2 \end{bmatrix},
\]
which recovers $\bA$, as required.

\section*{Problem 2.11}

Let $\bA$ be $p\times p$ diagonal, with $a_{ij}=0$ for $i\neq j$. We show $|\bA| = a_{11}a_{22}\cdots a_{pp}$ by induction on $p$.

\paragraph{Base case ($p=1$):} $\bA = [a_{11}]$, so $|\bA| = a_{11}$, trivially true.

\paragraph{Inductive step:} Assume the result holds for all diagonal matrices of size $(p-1)\times(p-1)$. By cofactor expansion along the first row (Definition 2A.24),
\[
|\bA| = a_{11}A_{11} + a_{12}A_{12} + \cdots + a_{1p}A_{1p}.
\]
Since $\bA$ is diagonal, $a_{12}=a_{13}=\cdots=a_{1p}=0$, so
\[
|\bA| = a_{11}A_{11},
\]
where $A_{11} = (-1)^{1+1}M_{11} = M_{11}$, and $M_{11}$ is the determinant of the submatrix obtained by deleting the first row and first column of $\bA$. This submatrix is itself a $(p-1)\times(p-1)$ diagonal matrix with diagonal entries $a_{22}, a_{33},\dots,a_{pp}$. By the induction hypothesis,
\[
M_{11} = a_{22}a_{33}\cdots a_{pp}.
\]
Therefore,
\[
|\bA| = a_{11}(a_{22}a_{33}\cdots a_{pp}) = a_{11}a_{22}\cdots a_{pp}.
\]
By induction, this holds for all $p \geq 1$. $\blacksquare$

\section*{Problem 2.24}

Given
\[
\bSigma = \begin{bmatrix} 4 & 0 & 0 \\ 0 & 9 & 0 \\ 0 & 0 & 1 \end{bmatrix}.
\]

\paragraph{(a) $\bSigma^{-1}$}
Since $\bSigma$ is diagonal,
\[
\bSigma^{-1} = \begin{bmatrix} 1/4 & 0 & 0 \\ 0 & 1/9 & 0 \\ 0 & 0 & 1 \end{bmatrix}.
\]

\paragraph{(b) Eigenvalues and eigenvectors of $\bSigma$}
Because $\bSigma$ is diagonal, its eigenvalues are the diagonal entries, with the standard basis vectors as corresponding eigenvectors. Ordered from largest to smallest:
\[
\lambda_1 = 9,\ \ e_1 = (0,1,0)'; \qquad
\lambda_2 = 4,\ \ e_2 = (1,0,0)'; \qquad
\lambda_3 = 1,\ \ e_3 = (0,0,1)'.
\]

\paragraph{(c) Eigenvalues and eigenvectors of $\bSigma^{-1}$}
$\bSigma^{-1}$ shares the same eigenvectors as $\bSigma$, with eigenvalues $1/\lambda_i$. Ordered from largest to smallest:
\[
\lambda_1^{*} = 1,\ \ e_1 = (0,0,1)'; \qquad
\lambda_2^{*} = \tfrac14,\ \ e_2 = (1,0,0)'; \qquad
\lambda_3^{*} = \tfrac19,\ \ e_3 = (0,1,0)'.
\]

\section*{Problem 2.30}

Given $\bmu_X' = [4,3,2,1]$,
\[
\bSigma_X = \begin{bmatrix} 3 & 0 & 2 & 2 \\ 0 & 1 & 1 & 0 \\ 2 & 1 & 9 & -2 \\ 2 & 0 & -2 & 4 \end{bmatrix},
\qquad
\bX^{(1)} = \begin{bmatrix} X_1\\X_2 \end{bmatrix}, \quad
\bX^{(2)} = \begin{bmatrix} X_3\\X_4 \end{bmatrix},
\]
\[
\bA = [1 \ \ 2], \qquad \bB = \begin{bmatrix} 1 & -2 \\ 2 & -1 \end{bmatrix}.
\]

\paragraph{(a) $E(\bX^{(1)})$}
\[
E(\bX^{(1)}) = \begin{bmatrix} 4 \\ 3 \end{bmatrix}.
\]

\paragraph{(b) $E(\bA\bX^{(1)})$}
\[
E(\bA\bX^{(1)}) = \bA\,E(\bX^{(1)}) = [1\ \ 2]\begin{bmatrix} 4\\3 \end{bmatrix} = 4 + 6 = 10.
\]

\paragraph{(c) $\mathrm{Cov}(\bX^{(1)})$}
Taking the rows/columns corresponding to $X_1, X_2$ from $\bSigma_X$:
\[
\mathrm{Cov}(\bX^{(1)}) = \begin{bmatrix} 3 & 0 \\ 0 & 1 \end{bmatrix}.
\]

\paragraph{(d) $\mathrm{Cov}(\bA\bX^{(1)})$}
\[
\mathrm{Cov}(\bA\bX^{(1)}) = \bA\,\mathrm{Cov}(\bX^{(1)})\,\bA'
= [1\ \ 2]\begin{bmatrix} 3&0\\0&1 \end{bmatrix}\begin{bmatrix} 1\\2 \end{bmatrix}
= [3\ \ 2]\begin{bmatrix} 1\\2 \end{bmatrix} = 3+4 = 7.
\]

\paragraph{(e) $E(\bX^{(2)})$}
\[
E(\bX^{(2)}) = \begin{bmatrix} 2 \\ 1 \end{bmatrix}.
\]

\paragraph{(f) $E(\bB\bX^{(2)})$}
\[
E(\bB\bX^{(2)}) = \bB\,E(\bX^{(2)}) = \begin{bmatrix} 1 & -2 \\ 2 & -1 \end{bmatrix}\begin{bmatrix} 2\\1 \end{bmatrix}
= \begin{bmatrix} 2-2 \\ 4-1 \end{bmatrix} = \begin{bmatrix} 0 \\ 3 \end{bmatrix}.
\]

\paragraph{(g) $\mathrm{Cov}(\bX^{(2)})$}
Taking the rows/columns corresponding to $X_3, X_4$ from $\bSigma_X$:
\[
\mathrm{Cov}(\bX^{(2)}) = \begin{bmatrix} 9 & -2 \\ -2 & 4 \end{bmatrix}.
\]

\paragraph{(h) $\mathrm{Cov}(\bB\bX^{(2)})$}
\[
\mathrm{Cov}(\bB\bX^{(2)}) = \bB\,\mathrm{Cov}(\bX^{(2)})\,\bB'.
\]
First,
\[
\bB\,\mathrm{Cov}(\bX^{(2)}) = \begin{bmatrix} 1 & -2 \\ 2 & -1 \end{bmatrix}\begin{bmatrix} 9 & -2 \\ -2 & 4 \end{bmatrix}
= \begin{bmatrix} 13 & -10 \\ 20 & -8 \end{bmatrix}.
\]
Then,
\[
\mathrm{Cov}(\bB\bX^{(2)}) = \begin{bmatrix} 13 & -10 \\ 20 & -8 \end{bmatrix}\begin{bmatrix} 1 & 2 \\ -2 & -1 \end{bmatrix}
= \begin{bmatrix} 33 & 36 \\ 36 & 48 \end{bmatrix}.
\]

\paragraph{(i) $\mathrm{Cov}(\bX^{(1)},\bX^{(2)})$}
Taking rows $1,2$ and columns $3,4$ of $\bSigma_X$:
\[
\mathrm{Cov}(\bX^{(1)},\bX^{(2)}) = \begin{bmatrix} 2 & 2 \\ 1 & 0 \end{bmatrix}.
\]

\paragraph{(j) $\mathrm{Cov}(\bA\bX^{(1)},\bB\bX^{(2)})$}
\[
\mathrm{Cov}(\bA\bX^{(1)},\bB\bX^{(2)}) = \bA\,\mathrm{Cov}(\bX^{(1)},\bX^{(2)})\,\bB'.
\]
First,
\[
\bA\,\mathrm{Cov}(\bX^{(1)},\bX^{(2)}) = [1\ \ 2]\begin{bmatrix} 2 & 2 \\ 1 & 0 \end{bmatrix} = [4\ \ 2].
\]
Then,
\[
[4\ \ 2]\begin{bmatrix} 1 & 2 \\ -2 & -1 \end{bmatrix} = [4-4 \ \ 8-2] = [0 \ \ 6].
\]
So
\[
\mathrm{Cov}(\bA\bX^{(1)},\bB\bX^{(2)}) = [0\ \ 6].
\]

\section*{Problem 2.41}

Given $\bmu_X' = [3,2,-2,0]$, $\bSigma_X = 3\bI_4$, and
\[
\bA = \begin{bmatrix} 1 & -1 & 0 & 0 \\ 1 & 1 & -2 & 0 \\ 1 & 1 & 1 & -3 \end{bmatrix}.
\]

\paragraph{(a) $E(\bA\bX)$}
\[
E(\bA\bX) = \bA\bmu_X =
\begin{bmatrix}
1(3)+(-1)(2)+0(-2)+0(0) \\
1(3)+1(2)+(-2)(-2)+0(0) \\
1(3)+1(2)+1(-2)+(-3)(0)
\end{bmatrix}
=
\begin{bmatrix} 1 \\ 9 \\ 3 \end{bmatrix}.
\]

\paragraph{(b) $\mathrm{Cov}(\bA\bX)$}
Since $\bSigma_X = 3\bI$,
\[
\mathrm{Cov}(\bA\bX) = \bA\bSigma_X\bA' = 3\,\bA\bA'.
\]
Computing $\bA\bA'$ entrywise (dot products of the rows of $\bA$):
\[
\bA\bA' = \begin{bmatrix} 2 & 0 & 0 \\ 0 & 6 & 0 \\ 0 & 0 & 12 \end{bmatrix}.
\]
Therefore,
\[
\mathrm{Cov}(\bA\bX) = 3\begin{bmatrix} 2 & 0 & 0 \\ 0 & 6 & 0 \\ 0 & 0 & 12 \end{bmatrix}
= \begin{bmatrix} 6 & 0 & 0 \\ 0 & 18 & 0 \\ 0 & 0 & 36 \end{bmatrix}.
\]

\paragraph{(c) Which pairs have zero covariance?}
Let $Y_1, Y_2, Y_3$ denote the three linear combinations given by the rows of $\bA$ (so $\bA\bX = (Y_1,Y_2,Y_3)'$). Since $\mathrm{Cov}(\bA\bX)$ is diagonal, every off-diagonal covariance is zero. Hence \textbf{all three pairs} $(Y_1,Y_2)$, $(Y_1,Y_3)$, and $(Y_2,Y_3)$ have zero covariance -- the three linear combinations are mutually uncorrelated.

\end{document}
